\documentclass[a4paper,10pt,exos]{nsi} 
\pagestyle{empty}


%%%%%%%%%%%%%%%%%%%%%%%%%%%%%%%%%%%%%%%%%%%%%%%%%
% Olympiades nationnales 2023 - Exercice 3 (candidats non spé)
%%%%%%%%%%%%%%%%%%%%%%%%%%%%%%%%%%%%%%%%%%%%%%%%%

\begin{document}
    \classe{Olympiades}
    \titre{Codes détecteurs et correcteurs}
    \maketitle

    {\titlefont\color{UGLiBlue}{Question préliminaire}}\\
    \question \textcolor{UGLiBlue}{Soient $a$ et $b$ deux nombres entiers.\\
    Montrer que le nombre $a+b$ est pair si, et seulement si, $a$ et $b$ sont de la même parité.}\\
    \textbf{Par disjonction de cas :}
    \begin{enumerate}[label=\textbullet]
        \item Si $a$ et $b$ sont pairs tous les deux :\\
        Il existe deux entiers $n$ et $p$ tels que $a=2n$ et $b=2p.$\\
        On a donc : $a+b=2(n+p)$, donc $a+b$ est pair.
        \item Si $a$ et $b$ sont impairs tous les deux :\\
        Il existe deux entiers $n$ et $p$ tels que $a=2n+1$ et $b=2p+1.$\\
        On a donc : $a+b=2(n+p) +2=2(n+p+1)$, donc $a+b$ est pair.
        \item Si $a$ et $b$ sont de parités différentes :\\
        Considérons par exemple le cas $a$ pair et $b$ impair : Il existe deux entiers $n$ et $p$ tels que $a=2n$ et $b=2p+1$.\\
        On a donc : $a+b=2(n+p)+1$, donc $a+b$ est impair.
    \end{enumerate}
    Ainsi $a+b$ est pair si, et seulement si, $a$ et $b$ sont de même parité.\\

    {\titlefont\color{UGLiBlue}{Codage d'un message}}\\
    \dleft{12cm}{
        \textcolor{UGLiBlue}{Un message est ici un nombre $M$ codé sous la forme d’un quadruplet $(x_1,x_2,x_3,x_4)$ où $x_1,x_2,x_3$ et $x_4$ sont des « bits », c'est-à-dire des nombres ne pouvant valoir que 0 ou 1. Le nombre $M$ que représente le quadruplet $(x_1,x_2,x_3,x_4)$, appelé aussi demi-octet d’information, vaut par définition :}
    }
    {\includegraphics[width=4.5cm]{binary.jpg}}
    $$  \textcolor{UGLiBlue}{M=x_1+2\times x_2+4\times x_3+8\times x_4.}$$
    \textcolor{UGLiBlue}{Par exemple, le code $(0,0,1,1)$ représente le nombre $M = 12$ puisque $12=0+2\times 0+4\times 1+8\times 1$.}\\

    \question \begin{enumalph}
        \item  \textcolor{UGLiBlue}{Quel est le message $M$ que code le quadruplet $(1,0,0,1)$ ?}\\
        Le quadruplet $(1,0,0,1)$ code le message $M=1+2\times 0+4\times 0+8\times 1=9$
        \item \textcolor{UGLiBlue}{Trouver un code qui représente $M = 10$. Trouver un code qui représente $M = 15$}.\\
        $10=2+8=0+2\times 1+4\times 0+8\times 1$. Le quadruplet $(0,1,0,1)$ représente $M=10$.\\
        $15=1+2+4+8=1+2\times 1+4\times 1+8\times 1$. Le quadruplet $(1,1,1,1)$ représente $M=15$.
        \item \textcolor{UGLiBlue}{Peut-on trouver un code pour représenter $M = 20$ ?}
        Le plus grand nombre que l'on peut représenter avec 4 un quadruplet est $1+2\times 1+4\times 1+8\times 1=15$. On ne peut pas représenter $M=20$.
        \item \textcolor{UGLiBlue}{Quels sont les différents messages possibles ?}\\
        Il y a 16 messages possibles : les 16 entiers de 0 à 15.
    \end{enumalph}

    \textcolor{UGLiBlue}{\small{\textit{Un message est parfois altéré (on dit aussi « corrompu ») lors de sa transmission du fait d’un matériel défectueux ou de signaux parasites. Des erreurs modifient des bits, un 0 se transformant en 1 ou un 1 se transformant en 0. Aussi des techniques permettant de détecter et de corriger ces anomalies ont-elles été mises au point. Ceci fait l’objet de la suite.}}}\\

    \normalsize

    {\titlefont\color{UGLiBlue}{Codage d'un message avec protection contre les erreurs}}\\
    \question \textbf{\textcolor{UGLiBlue}{Principe du bit de parité}}\\
    \textcolor{UGLiBlue}{Le code $(x_1,x_2,x_3,x_4)$ est transformé en le quintuplet $(x_1,x_2,x_3,x_4,y)$, dont le dernier bit $y$, dit de parité, vaut 0 si la somme $x_1+x_2+x_3+x_4$ est paire, et 1 si elle est impaire. C’est ce quintuplet qui est transmis, il représente le même message $M$ que le code $(x_1,x_2,x_3,x_4)$, à savoir $M=x_1+2\times x_2+4\times x_3+8\times x_4$. Les bits dits d’information demeurent $x_1,x_2,x_3,x_4$ et le bit de parité, $y$, est transmis avec les plus grandes précautions.\\[.5em]
    Par exemple, pour transmettre le nombre $M = 12$ correspondant à $x_1 = 0, x_2 = 0, x_3 = 1$ et $x_4 = 1$, on calcule d’abord $x_1 + x_2 + x_3 + x_4 = 2$, qui est pair ; on pose donc $y = 0$ et on émet le quintuplet $(0,0,1,1,0)$.}\\ 

    \question \textbf{\textcolor{UGLiBlue}{Principe des bits de contrôle}}\\
    \textcolor{UGLiBlue}{Le code $(x_1, x_2, x_3, x_4)$ est transformé en l’heptuplet $(x_1, x_2, x_3, x_4, y_1, y_2, y_3)$, où $y_1 = 0$ si $x_1 + x_2 + x_3$ est pair, $y_1 = 1$ sinon ; $y_2 = 0$ si $x_2 + x_3 + x_4$ est pair, $y_2 = 1$ sinon ; $y_3 = 0$ si $x_1 + x_3 + x_4$ est pair, $y_3 = 1$ sinon. Les bits dits d’information demeurent $x_1, x_2, x_3, x_4$.\\
    L’heptuplet $(x_1, x_2, x_3, x_4, y_1, y_2, y_3)$ code toujours le message $M = x_1 + 2 \times x_2 + 4 \times x_3 + 8 \times x_4$.}

    \begin{enumalph}
        \item \textcolor{UGLiBlue}{Quels sont les bits $y_1, y_2, y_3$, dits de contrôle, associés au quadruplet $(1,0,0,1)$ codant le nombre $M = 9$ ?}\\
        On a : 
        \begin{enumerate}[label=\textbullet]
            \item $x_1+x_2+x_3=1+0+0=1$, donc $y_1=1$.
            \item $x_2+x_3+x_4=0+0+1=1$, donc $y_2=1$.
            \item $x_1+x_3+x_4=1+0+1=2$, donc $y_3=0$.
        \end{enumerate}
        \item \textcolor{UGLiBlue}{Pourquoi est-on certain que l’heptuplet reçu $(1,1,0,1,0,0,1)$ résulte d’une altération de transmission dans le cas où on est sûr des bits de contrôle ?}\\
        Ici $x_1+x_3+x_4=1+0+1=2$, donc $y_3$ devrait être égal à $0$. Ce n'est pas le cas dans le triplet transmis.
        \item \textcolor{UGLiBlue}{Si on est sûr de la justesse des bits de contrôle, dans l’hypothèse où exactement un des quatre bits d’information est erroné, pourquoi peut-on détecter qu’il y a eu une altération et pourquoi peut-on la localiser (et donc la corriger) ? Peut-on détecter l’erreur quand exactement deux des quatre bits d’information sont erronés ?}\\
        Dans le cas où exactement un des quatre bits d'information est erroné :
        \begin{enumerate}[label=\textbullet]
            \item Si $x_1$ est erroné, alors $y_1$ et $y_3$ ne correspondront pas au quadruplet $(x_1,x_2,x_3,x_4)$.
            \item Si $x_2$ est erroné, alors $y_1$ et $y_2$ ne correspondront pas au quadruplet $(x_1,x_2,x_3,x_4)$.
            \item Si $x_3$ est erroné, alors $y_1,y_2$ et $y_3$ ne correspondront pas au quadruplet $(x_1,x_2,x_3,x_4)$.
            \item Si $x_4$ est erroné, alors $y_2$ et $y_3$ ne correspondront pas au quadruplet $(x_1,x_2,x_3,x_4)$.
        \end{enumerate}
        On peut donc localiser l'erreur à partir du calcul de $y_1, y_2$ et $y_3$.\\
        
        Dans le cas où exactement deux des quatre bits d'information sont erronés :
        \begin{enumerate}[label=\textbullet]
            \item Si $x_1$ et $x_2$ sont erronés, alors $y_2$ et $y_3$ ne correspondront pas au quadruplet $(x_1,x_2,x_3,x_4)$.
            \item Si $x_1$ et $x_3$ sont erronés, alors $y_2$ ne correspondra pas au quadruplet $(x_1,x_2,x_3,x_4)$.
            \item Si $x_1$ et $x_4$ sont erronés, alors $y_1$ et $y_2$ ne correspondront pas au quadruplet $(x_1,x_2,x_3,x_4)$.
            \item Si $x_2$ et $x_3$ sont erronés, alors $y_3$ ne correspondra pas au quadruplet $(x_1,x_2,x_3,x_4)$.
            \item Si $x_2$ et $x_4$ sont erronés, alors $y_1$ et $y_3$ ne correspondront pas au quadruplet $(x_1,x_2,x_3,x_4)$.
            \item Si $x_3$ et $x_4$ sont erronés, alors $y_1$ ne correspondra pas au quadruplet $(x_1,x_2,x_3,x_4)$.
        \end{enumerate}
        On peut donc localiser l'erreur à partir du calcul de $y_1, y_2$ et $y_3$.
    \end{enumalph}



\end{document}